% ============================================================================
%  宏包集合：preamble/packages.tex
%  原则：① 只用 TeX Live 主流发行版默认就有的宏包；
%        ② 一个宏包只加载一次，避免重复 \usepackage；
%        ③ 冲突宏包（如 caption/subcaption 与部分期刊类）集中在此处便于排查。
%  编译引擎：XeLaTeX（中文必需，故不加载 inputenc / fontenc）
% ============================================================================

% --- 工具宏（\AtBeginEnvironment 等）---
\usepackage{etoolbox}

% --- 排版基础 ---
\usepackage{xcolor}               % 颜色（其余宏包依赖它）
\usepackage{microtype}            % 字符 protrusion，改善中文排版灰度
\usepackage{geometry}             % 页边距（layout.tex 中设置）

% --- 数学 ---
\usepackage{amsmath, amssymb}     % 公式与数学符号
\usepackage{amsthm}               % 定理/证明/引理环境
\usepackage{mathtools}            % amsmath 增强：\coloneqq, \DeclarePairedDelimiter 等
\usepackage{bm}                   % 加粗数学符号（向量/矩阵）

% --- 插图与表格 ---
\usepackage{graphicx}             % 插图
\usepackage{subcaption}           % 子图（若与期刊模板冲突，改用 \subfigure 或并列 minipage）
\usepackage{booktabs}             % 三线表
\usepackage{array}                % 列格式控制
\usepackage{multirow}             % 跨行单元格
\usepackage{makecell}             % 单元格内换行 \makecell{}
\usepackage{threeparttable}       % 表注（置于表格宽度内）
\usepackage{longtable}            % 跨页长表（符号表适用）
\usepackage{tabularx}             % 自适应宽度列

% --- 算法伪代码 ---
\usepackage{algorithm}
\usepackage{algpseudocode}

% --- 单位与数值 ---
\usepackage{siunitx}              % \SI{}{\metre\per\second}、\num{}
\sisetup{
  range-phrase = {$\sim$},
  list-final-separator = { 和 },
  output-decimal-marker = {.}
}

% --- 中文支持 ---
\usepackage[UTF8]{ctex}           % XeLaTeX + UTF-8 源码，勿改

% --- 列表 ---
\usepackage{enumitem}
\setlist[enumerate]{labelsep=0.6em, leftmargin=2.2em}
\setlist[itemize]{labelsep=0.6em, leftmargin=2.2em}

% --- 审稿辅助（draft 模式下由 draft-tools.tex 决定是否启用）---
\ifdraft
  \usepackage{lineno}
  \modulolinenumbers[5]
\fi
